% =====================================================================
%  Préambule commun (modèle) : les données du chapitre sont dans chapitre.tex
%  Chapitre de manuel de Physique-Chimie (1re année du baccalauréat)
%  Préambule commun au livre du professeur et au livre de l'élève :
%  \def\versioneleve{1} en première ligne active la version élève (à compléter).
%  Compilation : ./build_livre.sh (XeLaTeX, deux passes)
%  Les figures viennent des scènes animées de la présentation
%  (figures_html/book.html + render_book.py) et sont exportées dans images/.
% =====================================================================
\documentclass[11pt,a4paper]{article}
% Données du chapitre (lues par preamble.tex)
\newcommand{\MatiereLivre}{PHYSIQUE}
\newcommand{\DomaineLivre}{Mécanique}
\newcommand{\NumChapLivre}{3}
\newcommand{\TitreCourtLivre}{Travail et énergie cinétique}
% \MatiereLivre, \DomaineLivre, \NumChapLivre, \TitreCourtLivre

% ---------- Mise en page ----------
\usepackage[a4paper,top=2.3cm,bottom=2.2cm,left=1.9cm,right=2.1cm,headsep=6mm,footskip=11mm]{geometry}
\usepackage{fontspec}
\usepackage{polyglossia}
\setmainlanguage{french}
\usepackage{amsmath}
\usepackage{unicode-math}
\setmainfont{TeX Gyre Pagella}
\setsansfont{Poppins}[Path=fonts/, Extension=.ttf, UprightFont=*-Regular, BoldFont=*-Bold, ItalicFont=*-Italic, BoldItalicFont=*-BoldItalic]
\newfontfamily\titlefont{Poppins}[Path=fonts/, Extension=.ttf, UprightFont=*-Bold, BoldFont=*-Bold]
\setmathfont{Latin Modern Math}
\setmathfont[range={it/{latin,Latin,greek,Greek}}]{texgyrepagella-italic.otf}
\setmathfont[range={up/{latin,Latin,num}}]{texgyrepagella-regular.otf}

\usepackage[table]{xcolor}
\usepackage{graphicx}
\graphicspath{{images/}}
\usepackage{tikz}
\usetikzlibrary{calc,positioning,shapes.misc}
\usepackage[most]{tcolorbox}
\usepackage{enumitem}
\usepackage{fancyhdr}
\usepackage{titlesec}
\usepackage{eso-pic}
\usepackage{multicol}
\usepackage{array,tabularx,booktabs}
\usepackage{needspace}
\usepackage{ragged2e}
\setlength{\RaggedRightRightskip}{0pt plus 6em}
\usepackage{microtype}
\usepackage{siunitx}
\sisetup{locale=FR, per-mode=power, inter-unit-product=\ensuremath{\mathord{\cdot}}, group-minimum-digits=4}
\usepackage[hidelinks]{hyperref}

\setlength{\parindent}{0pt}
\setlength{\parskip}{3pt plus 1pt}
\linespread{1.07}
\emergencystretch=3em
\tolerance=2500

% ---------- Couleurs (identiques à la charte des figures) ----------
\definecolor{cInk}{HTML}{1B2A41}
\definecolor{cInk2}{HTML}{51627A}
\definecolor{cBlue}{HTML}{1F5FA8}
\definecolor{cSky}{HTML}{8EC5FF}
\definecolor{cOrange}{HTML}{F07F1E}
\definecolor{cRed}{HTML}{D62839}
\definecolor{cGreen}{HTML}{1F9D74}
\definecolor{cPurple}{HTML}{6C4AB6}
\definecolor{cTeal}{HTML}{0E8F8A}
\definecolor{cAmber}{HTML}{F4A300}
\definecolor{cCyan}{HTML}{0B84C6}
\definecolor{cGrey}{HTML}{8A96A8}
\color{cInk}

% ---------- Petites commandes ----------
\newcommand\og{«\nobreak\thinspace}
\newcommand\fg{\nobreak\thinspace»}
\newcommand{\arc}[1]{\overparen{#1}}
\newcommand{\vect}[1]{\overrightarrow{#1}}
\newcommand{\Mom}{\mathcal{M}}
\newcommand{\Ec}{E_{\mathrm{C}}}
\newcommand{\kw}[1]{{\sffamily\bfseries\color{cBlue}#1}}   % mot-clé des questions
\newcommand{\hl}[1]{{\bfseries\color{cRed}#1}}             % mise en évidence
\newcommand{\imp}[1]{\textbf{#1}}
\newcommand{\sfb}[1]{{\sffamily\bfseries #1}}

% Pastille « Doc n »
\newcommand{\docbadge}[2][cBlue]{\tikz[baseline=(b.base)]\node[fill=#1,text=white,rounded corners=6pt,inner xsep=5pt,inner ysep=2.2pt,font=\sffamily\bfseries\footnotesize](b){Doc #2};}
% Document illustré : [largeur]{n}{légende}{image}
\newcommand{\docimg}[4][0.485\linewidth]{%
  \begin{minipage}[t]{#1}\centering
    \includegraphics[width=\linewidth]{#4}\par\vspace{1.4mm}
    {\sffamily\small\docbadge{#2}\enspace\color{cInk2}#3}
  \end{minipage}}
% Figure numérotée simple
\newcommand{\figcap}[1]{\par\vspace{1mm}{\centering\sffamily\small\color{cInk2}#1\par}}

% Formule encadrée
\newtcbox{\formule}[1][cRed]{on line, enhanced, colback=white, colframe=#1, boxrule=1pt, arc=2.5mm,
  left=2.5mm,right=2.5mm,top=1.2mm,bottom=1.2mm, boxsep=0pt}

% ---------- Rubriques (boîtes) ----------
% \begin{rub}[options]{couleur}{icône}{titre}
\newtcolorbox{rub}[4][]{%
  enhanced, breakable,
  colback=#2!4!white, colframe=#2!75!white, boxrule=0.9pt, arc=3mm,
  left=4.5mm, right=4.5mm, top=6.5mm, bottom=3.5mm,
  before skip=6mm, after skip=4mm,
  title={\raisebox{-2.35mm}{\includegraphics[height=7.6mm]{ic_#3}}\hspace{2.2mm}#4},
  fonttitle=\sffamily\bfseries\fontsize{11.5}{13}\selectfont\color{white},
  attach boxed title to top left={xshift=5mm,yshift=-4.4mm},
  boxed title style={enhanced, colback=#2, colframe=#2, arc=4.3mm, boxrule=0pt,
     left=0.4mm, right=4.5mm, top=0.5mm, bottom=0.5mm},
  #1}

% Encadré secondaire (méthode, rappel, le saviez-vous) — non sécable
\newtcolorbox{encart}[4][]{%
  enhanced, colback=#2!6!white, colframe=#2!6!white, boxrule=0pt, arc=2.5mm,
  left=13mm, right=3mm, top=2mm, bottom=2mm, before skip=3mm, after skip=3mm,
  borderline west={1.6pt}{0pt}{#2},
  overlay={\node[anchor=north west] at ([xshift=2mm,yshift=-2mm]frame.north west)
     {\includegraphics[width=8.5mm]{ic_#3}};},
  fontupper=\small,
  title={#4}, detach title, coltitle=#2, before upper={{\sffamily\bfseries\color{#2}\tcbtitle}\par\smallskip},
  #1}

% Situation de départ (bulle)
\newtcolorbox{situation}{%
  enhanced, colback=cAmber!9!white, colframe=cAmber!9!white, boxrule=0pt, arc=3mm,
  left=16mm, right=4mm, top=2.5mm, bottom=2.5mm, before skip=2mm, after skip=3mm,
  overlay={\node at ([xshift=8mm]frame.west) {\includegraphics[width=10mm]{ic_situation}};
           \draw[cAmber, line width=1.4pt, rounded corners=3mm] (frame.south west) rectangle (frame.north east);}}

% Sous-titre à l'intérieur d'une rubrique
\newcommand{\soustitre}[2][cTeal]{\par\medskip{\sffamily\bfseries\color{#1}\MakeUppercase{#2}}\par\smallskip}

% Listes
% pastille numérotée (commande robuste : enumitem développe l'étiquette)
\DeclareRobustCommand{\pastille}[2]{\tikz[baseline=(n.base)]\node[circle,fill=#1,text=white,inner sep=0pt,minimum size=5.2mm,font=\sffamily\bfseries\footnotesize](n){#2};}
\setlist[enumerate,1]{label={\pastille{cBlue}{\arabic*}},
  leftmargin=8mm, labelsep=2.5mm, itemsep=3pt, topsep=3pt}
\newlist{reps}{enumerate}{1}
\setlist[reps]{label={\pastille{cGreen}{\arabic*}},
  leftmargin=8mm, labelsep=2.5mm, itemsep=4pt, topsep=3pt}
\newlist{appq}{enumerate}{1}
\setlist[appq]{label={\sffamily\bfseries\color{cPurple}\arabic*.}, leftmargin=6mm, itemsep=2pt, topsep=3pt}
\newlist{exq}{enumerate}{1}
\setlist[exq]{label={\sffamily\bfseries\color{cOrange}\arabic*.}, leftmargin=6mm, itemsep=2pt, topsep=3pt}
\setlist[itemize,1]{label={\color{cOrange}\small$\blacktriangleright$}, leftmargin=6mm, itemsep=2pt, topsep=2pt}
\setlist[itemize,2]{label={\color{cInk2}\small$\bullet$}, leftmargin=5mm, itemsep=1pt, topsep=1pt}

% ---------- Titres de sections ----------
\titleformat{\section}[block]
  {\sffamily\bfseries\fontsize{18}{22}\selectfont\color{cInk}\RaggedRight}
  {\tikz[baseline=(n.base)]\node[fill=cOrange,rounded corners=3pt,text=white,minimum width=10.5mm,minimum height=10.5mm,inner sep=0pt,font=\sffamily\bfseries\fontsize{18}{22}\selectfont](n){\thesection};}
  {4mm}{}[\vspace{1.2mm}{\color{cOrange}\titlerule[1.3pt]}]
\titlespacing*{\section}{0pt}{8mm}{4mm}

% ---------- En-têtes et pieds de page ----------
\pagestyle{fancy}
\fancyhf{}
\fancyhead[L]{\sffamily\footnotesize\color{cInk2}\textcolor{cOrange}{\textbf{\MatiereLivre}}\enspace·\enspace \DomaineLivre}
\fancyhead[R]{\sffamily\footnotesize\color{cInk2}Chapitre \NumChapLivre\enspace·\enspace \TitreCourtLivre}
\renewcommand{\headrulewidth}{0.6pt}
\renewcommand{\headrule}{{\color{cBlue!25}\hrule height\headrulewidth width\headwidth}}
\fancyfoot[C]{\tikz\node[circle,fill=cBlue,text=white,font=\sffamily\bfseries\small,minimum size=7.5mm,inner sep=0pt]{\thepage};}
% onglet latéral
\AddToShipoutPictureBG{%
  \ifnum\value{page}>1
  \begin{tikzpicture}[remember picture,overlay]
    \fill[cOrange,rounded corners=2.5mm] ([xshift=-6.5mm,yshift=-3.2cm]current page.north east) rectangle ([xshift=3mm,yshift=-7.2cm]current page.north east);
    \node[rotate=-90,text=white,font=\sffamily\bfseries\footnotesize] at ([xshift=-3.3mm,yshift=-5.2cm]current page.north east) {CHAPITRE \NumChapLivre};
  \end{tikzpicture}\fi}


% ---------- Encadré « Rappel » ----------
\newenvironment{rappel}{\begin{encart}{cGrey}{rappel}{Rappel}}{\end{encart}}

% ---------- Tableaux de mesures : \mesures[largeur]{nb colonnes}{lignes} ----------
\newcommand{\mesures}[3][1.25cm]{{\renewcommand{\arraystretch}{1.35}\setlength{\arrayrulewidth}{.6pt}\arrayrulecolor{cTeal!55}\setlength{\tabcolsep}{4pt}%
  \begin{tabular}{|>{\sffamily\footnotesize\bfseries\color{cTeal}}l|*{#2}{>{\centering\arraybackslash}p{#1}|}}\hline #3\end{tabular}}}

% =====================================================================
%  Version élève (livre à compléter) : \def\versioneleve{1} avant \documentclass
% =====================================================================
\newif\ifeleve
\ifdefined\versioneleve\elevetrue\fi
\newlength{\interligne}\setlength{\interligne}{7mm}
\newsavebox{\repbox}
\newlength{\replen}
\newcount\repn \newcount\lgn
\newcommand{\pointilles}{{\color{cInk2!60}\leaders\hbox to 1.8mm{\hss.\hss}\hfill}\kern0pt}
% \lignes{n} : n lignes pointillées (la première continue la ligne en cours)
\newcommand{\lignes}[1]{\leavevmode\vrule width0pt height\dimexpr\interligne-1mm\relax depth0pt\pointilles
  \lgn=1 \loop\ifnum\lgn<#1 \par\nointerlineskip\noindent\vrule width0pt height\interligne depth0pt\pointilles\advance\lgn by 1 \repeat}
% \rep{réponse} : réponse rédigée (professeur) ou lignes pointillées en nombre adapté (élève)
\newcommand{\rep}[1]{%
  \ifeleve
    \sbox{\repbox}{\parbox{\linewidth}{#1}}%
    \replen=\dimexpr\ht\repbox+\dp\repbox\relax
    \repn=\numexpr\replen/\baselineskip+1\relax
    \ifnum\repn<2 \repn=2\fi
    \edef\nblignes{\the\repn}\lignes{\nblignes}%
  \else #1\fi}
% \trou{mot ou formule} : mot clé des « À retenir » (professeur) ou espace pointillé (élève)
\newlength{\troumax}
\newcommand{\trou}[1]{%
  \ifeleve
    \sbox{\repbox}{#1}%
    \replen=\dimexpr\wd\repbox*3/2+6mm\relax
    \troumax=\linewidth \ifdim\hsize<\troumax \troumax=\hsize\fi
    \advance\troumax by -14mm\relax
    \ifdim\replen>\troumax \replen=\troumax\fi
    \makebox[\replen][l]{\vphantom{\usebox{\repbox}}\pointilles}%
  \else #1\fi}
% \troul{fin de phrase} : fin de phrase à compléter sur une ou plusieurs lignes
\newcommand{\troul}[1]{%
  \ifeleve
    \sbox{\repbox}{#1}%
    \repn=\numexpr\dimexpr\wd\repbox*3/2\relax/\linewidth+1\relax
    \edef\nblignes{\the\repn}\unskip\ \lignes{\nblignes}%
  \else #1\fi}
% \cel{valeur} : case de tableau (vide en version élève, hauteur conservée)
\newcommand{\cel}[1]{\ifeleve\vphantom{#1}\else #1\fi}
\newcommand{\titrereponses}{\ifeleve Mes réponses\else Réponses\fi}
\newcommand{\titrecorriges}{\ifeleve Mes réponses aux exercices\else Éléments de correction\fi}
\ifeleve
  \fancyhead[R]{\sffamily\footnotesize\color{cInk2}Chapitre \NumChapLivre\enspace·\enspace\textcolor{cTeal}{\textbf{Livre de l’élève}}}
\fi
